\subsection{A literal characteristic surface realization of the decisive modular quotient}

\begin{result}{A.4.10}{Proved}{An explicit finite four-puncture surface quotient retains the signed modular centralizer detector}
Let $B=C_2*C_3$, and let
\[
 D=B/(N_{12}'N_{12}^3),\qquad
 N_{12}=\ker\bigl(B\longrightarrow
                 \operatorname{PSL}_2(\mathbb Z/12)\bigr)
                                                        \tag{ASR.1}
\]
be the literal quotient of A.4.6, of order $576\cdot3^{97}$.
The projective linear group in (ASR.1) is the quotient by the
single diagonal subgroup $\{\pm I\}$.
Let $D^+$ be its signed pullback along
$\operatorname{SL}_2(\mathbb Z/12)\to
 \operatorname{PSL}_2(\mathbb Z/12)$; thus
$|D^+|=1152\cdot3^{97}$.

There is an explicitly specified characteristic finite quotient
$\mathcal Q$ of the free fundamental group $\Pi_{1,4}$ of
the torus with its four two-torsion points removed, and an
explicit surjection
\[
 J=\operatorname{im}\bigl(
 \operatorname{SL}_2(\mathbb Z)\longrightarrow
                 \operatorname{Out}(\mathcal Q)\bigr)
 \longrightarrow D^+.                                  \tag{ASR.2}
\]
The modular action is the actual two-division normalizer action,
with the origin puncture fixed.  Its generator images in (ASR.2)
are the specified signed modular generator images of A.4.6.
In particular this is a map from an actual finite surface
representation onto that modular quotient, not an identification
of their kernels without a morphism.
There is also an explicitly specified characteristic quotient
$\mathcal P$ of the once-punctured torus group itself, of order
$4|\mathcal Q|$, whose ordinary outer modular image maps onto
the same $D^+$.  It retains the actual two-, three-, and
six-division subgroups simultaneously, as described in
(ASR.17)--(ASR.20) below.

Here is a completely explicit index bound.  Put
\[
 d=96\cdot3^{97},\qquad
 e=d\,5^{d+1},\qquad
 r=e\,5^{1+2e}.
                                                        \tag{ASR.3}
\]
The construction gives
$|\mathcal Q|\leq(r!)^{(r!)^5}$ and
$|\mathcal P|\leq4(r!)^{(r!)^5}$.
It uses only the actual finite group (ASR.1), two specified
homology covers, and a specified finite diagonal map.
No termination assertion for an unspecified congruence search
enters the construction.
\end{result}

\begin{proof}
\par\medskip\noindent\textbf{The exact braid, free-group, and period maps.}
Let
$S=F(x_1,x_2,x_3)$ and $t=x_1x_2x_3$.
The braid generators $a,b$ act by
\[
 a:(x_1,x_2,x_3)\longmapsto
       (x_1x_2x_1^{-1},x_1,x_3),\qquad
 b:(x_1,x_2,x_3)\longmapsto
       (x_1,x_2x_3x_2^{-1},x_2).                        \tag{ASR.4}
\]
These are literal automorphisms, they satisfy $aba=bab$, and
they fix $t$.  The central braid $z=(ab)^3$ acts as
$\operatorname{Inn}(t)$.
Fill the fourth sphere puncture by the actual quotient
\[
 p_0:S\twoheadrightarrow T=F(x,y),\qquad
 (x_1,x_2,x_3)\longmapsto
                (x,y,(xy)^{-1}).                       \tag{ASR.5}
\]
On $T$ the two pure braid words have the exact actions
\[
 a^2=\operatorname{Inn}(xy),\qquad
 b^2=\operatorname{Inn}(x^{-1}).                        \tag{ASR.6}
\]
The full twist is trivial on $T$, so this is an action of
$B=B_3/\langle z\rangle$.
Its quotient on the three nonzero peripheral homology classes
modulo two is the usual $B\twoheadrightarrow\mathfrak S_3$.
Denote its kernel by $B(2)$.
The Bass--Serre cover has index six, trivial vertex stabilizers,
and Euler characteristic $-1$; hence $B(2)$ is free of rank two.
Equations (ASR.6) map it onto $\operatorname{Inn}(T)\cong T$,
because $xy,x^{-1}$ form a free basis.
A surjection between free groups of the same finite rank is an
isomorphism.  Thus (ASR.6) gives the actual equivariant
identification
$j:B(2)\xrightarrow{\sim}T$.
In the standard free-product marking $s=aba$, $r_0=ab$,
the braid words are $a=r_0^{-1}s$, $b=sr_0^2$.
These formulas specify their images in (ASR.1).

For the torus cover, take the parity kernel
\[
 \Pi=\ker\bigl(S\longrightarrow C_2,
                         \ x_i\longmapsto1\bigr).
                                                        \tag{ASR.7}
\]
This is the unbranched double cover of the four-punctured
sphere, which compactifies to the cover branched at the four
deleted points.  Its compactification has genus one and its
four deleted points are the fixed points of the elliptic
involution.  Choosing the fourth as origin identifies these
points with the full two-torsion subgroup.
The following is a literal free basis of $\Pi$:
\[
 d_1=x_1^2,\quad u=x_2x_1^{-1},\quad v=x_3x_1^{-1},
 \quad w=x_1x_2,\quad q=x_1x_3.                        \tag{ASR.8}
\]
Its four peripheral words are
$d_1,uw,vq,(wvd_1uq)^{-1}$.
After filling these punctures, $w=u^{-1}$, $q=v^{-1}$ and
the remaining relation is $u^{-1}vuv^{-1}=1$.
Use the primitive compact period basis $A=u$, $B=u^{-1}v$.
The exact period map, in the column basis (ASR.8), is
\[
 P=\begin{pmatrix}0&1&1&-1&-1\\0&0&1&0&-1\end{pmatrix}.
                                                        \tag{ASR.9}
\]
Direct restriction of (ASR.4) gives
$P a_*=U P$ and $P b_*=L P$, where
\[
 U=\begin{pmatrix}1&1\\0&1\end{pmatrix},\qquad
 L=\begin{pmatrix}1&0\\-1&1\end{pmatrix}.
                                                        \tag{ASR.10}
\]
The full twist acts by $-I$ on these periods, and its square
acts on $\Pi$ as the literal inner automorphism by $t^2$.
Consequently the outer action factors through
$B_3/\langle z^2\rangle=\operatorname{SL}_2(\mathbb Z)$,
with exactly the signed marking (ASR.10).
All free-word and period identities just used are checked in
\path{modular_artin_surface_bridge.py}; no change of the
ordinary homology marking is suppressed.

\par\medskip\noindent\textbf{Two explicit finite homology refinements.}
Let $q_D:B\twoheadrightarrow D$ be (ASR.1), and put
$D_0=q_D(B(2))$.  Its order is $d$ in (ASR.3), since
$D$ retains the displayed quotient to $\mathfrak S_3$.
The actual quotient of $T$ corresponding to this subgroup is
\[
 p:T\twoheadrightarrow D_0,\qquad
 p(x)=q_D(b^{-2}),\qquad p(y)=q_D(b^2a^2).             \tag{ASR.11}
\]
Indeed $j(b^{-2})=x$ and $j(b^2a^2)=y$ by (ASR.6).
Equivariance of $j$ shows that $N_T=\ker p$ is preserved
by the full $B$-action on $T$.
It is free of rank $d+1$.
Define the actual finite group and quotient
\[
 D_1=T/(N_T'N_T^5),\qquad p_1:T\twoheadrightarrow D_1.
                                                        \tag{ASR.12}
\]
Its order is $e=d5^{d+1}$; its map to $D_0$ and its
$B$-action are the induced maps, not independent choices.

The centre of $D_1$ maps trivially to $D_0$.
To prove this, use the $D_0$-cover of the two-loop rose.
Its first homology over $\mathbb F_5$ is the kernel
$N_T/(N_T'N_T^5)$ in (ASR.12).
For any nonidentity deck element, the vertex and edge
permutation traces are both zero.  The exact chain sequence
therefore gives trace one on first homology.  The identity
has trace $d+1$, which differs from one because $5\nmid d$.
The deck action on this homology is consequently faithful.
A central element of $D_1$ acts trivially on that kernel by
conjugation, so its image in $D_0$ must be the identity.

Now use the literal composite
$S\xrightarrow{p_0}T\xrightarrow{p_1}D_1$ and its kernel
$N_S$.  It has index $e$ in the rank-three free group and
therefore rank $1+2e$.  Define
\[
 R=S/(N_S'N_S^5),\qquad p_R:S\twoheadrightarrow R.
                                                        \tag{ASR.13}
\]
Its order is $r$ in (ASR.3).
The Artin action preserves its kernel, because it preserves
$t$ and induces the already specified action on $D_1$.

The exact centralizer assertion required here is
\[
 C_R(p_R(t))\longrightarrow D_1\quad\text{has image }1.
                                                        \tag{ASR.14}
\]
In the free basis $x_1,x_2,t$ of $S$, its $D_1$-cover has
one $t$-loop at every vertex.  Choose a spanning tree using
only $x_1,x_2$ edges; they connect the cover because
$p_1(T)=D_1$.
The $t$-loops are then independent members of the actual
Schreier homology basis.  The class of $t$ is the loop at
the identity vertex.  Conjugating it by an element with
residue $h\in D_1$ gives the loop at $h$.
Equality of these two classes forces $h=1$, proving
(ASR.14) directly in the finite group (ASR.13).

\par\medskip\noindent\textbf{The modular quotient detected by innerness.}
Suppose an Artin automorphism indexed by a braid is inner
on $R$.  This is an implication between explicit finite
group actions; we now compute the image of every such braid.
Write the inner automorphism as $\operatorname{Inn}(v)$.
It fixes $p_R(t)$, so (ASR.14) gives $v\mapsto1$ in $D_1$.
The induced automorphism of $D_1$ is therefore the identity.
The quotient $D_1\to D_0$ retains the actual quotient
$T\to H_1(T,\mathbb F_2)$: in modular terms this is the
image $B(2)/B(4)\cong C_2^2$ obtained by reduction modulo
four and then by the common sign.
The action of $B/B(2)=\mathfrak S_3$ on its three nonzero
vectors is faithful.  The braid is consequently pure
modulo the central full twist.

For its class $g\in B(2)$, the induced automorphism of
$D_1$ is exactly $\operatorname{Inn}(p_1(j(g)))$ by
(ASR.6).  Its being the identity says that
$p_1(j(g))\in Z(D_1)$.
The centre calculation after (ASR.12) gives
$p(j(g))=1$.  Equation (ASR.11) now gives $q_D(g)=1$.
Thus the kernel of the actual outer Artin action on $R$
is contained in the kernel of the specific map $q_D$.

\par\medskip\noindent\textbf{A characteristic genus-one surface quotient.}
The restriction $\Pi\to R$ is onto.
Indeed $t\in N_S$, so $t^5\in\ker p_R$, whereas $t^5$
has odd parity in (ASR.7).  Every coset of $\ker p_R$
therefore has an even representative.

Put $m=r!$.  List all $m^5$ assignments of five permutations
in $\mathfrak S_r$ to the five ordered free generators in
(ASR.8), and take the actual diagonal homomorphism
\[
 \Pi\longrightarrow(\mathfrak S_r)^{m^5}.
                                                        \tag{ASR.15}
\]
Define $\mathcal Q$ to be its image.
This fully specifies a finite group and its quotient map;
for example, the permutations can be ordered lexicographically.
Its kernel is characteristic, also under continuous
automorphisms of the full profinite free group: precomposition
permutes the complete finite list of homomorphisms.
Its order is at most $m^{m^5}$.
The surjection $\Pi\to R$, followed by the left regular
representation of $R$ on its $r$ elements, is one coordinate
of (ASR.15).  It supplies the actual quotient map
$\mathcal Q\twoheadrightarrow R$.

An element of $\operatorname{SL}_2(\mathbb Z)$ acting
innerly on $\mathcal Q$ acts innerly on $R$, by this map
and the previously verified invariance of its kernel.
The preceding paragraph proves that its projective image
in $D$ is the identity.  Its signed matrix modulo twelve
is therefore either $I$ or $-I$.
Since $r\geq3$, (ASR.15) includes every homomorphism
$\Pi\to C_3$ followed by a fixed embedding
$C_3\hookrightarrow\mathfrak S_r$.
Thus $\mathcal Q$ also retains
$H_1(\Pi,\mathbb F_3)$ and its actual compact-period
quotient (ASR.9).  An inner automorphism is the identity
on this homology.  Equations (ASR.9)--(ASR.10) therefore
exclude the signed value $-I$ modulo twelve.
The signed value in $D^+$ is exactly the identity.

We have proved kernel containment for the two literal
homomorphisms from $\operatorname{SL}_2(\mathbb Z)$.
It defines the surjection (ASR.2), with
\[
 \rho(a)\longmapsto(q_D(a),U\bmod12),\qquad
 \rho(b)\longmapsto(q_D(b),L\bmod12).                 \tag{ASR.16}
\]
Both targets are generated by these displayed images, so
this is also a finite algorithm for the map: enumerate the
finite generated image $J$ and transport each generator
word by (ASR.16).  Kernel containment proves independence
of the word.

\par\medskip\noindent\textbf{The ordinary once-punctured quotient and the actual common cover.}
The geometric curve underlying (ASR.7) has a specified origin,
namely the point over the fourth branch point.  In the
Weierstrass family
$Y^2=(X-e_1)(X-e_2)(X-e_3)$, with origin at infinity,
the other three branch points lift to exactly the three
nonzero two-torsion points.  Thus the Artin monodromy in
(ASR.4), lifted to this double cover, is the actual monodromy
of $E\setminus E[2]$.
Multiplication by two is a morphism of this same family;
its covering monodromy identifies that action with the
origin-preserving lift of the ordinary modular action.
The origin condition removes the possible deck translation,
because the four deck translations act freely on the four
deleted points.  This is the geometric comparison of the
two actions, in addition to the word and period checks
(ASR.8)--(ASR.10).

Let $\Sigma=\pi_1(E\setminus\{0\})$, with the based
covering maps fixed below.  The actual multiplication
cover $[2]$ identifies $\Pi$ with
$K_2=\ker(\Sigma\to H_1(E,\mathbb Z/2))$.
More generally
\[
 K_m=\ker\bigl(\Sigma\to H_1(E,\mathbb Z/m)\bigr)
   =[m]_*\pi_1(E\setminus E[m]),\qquad [\Sigma:K_m]=m^2.
                                                        \tag{ASR.17}
\]
This equality follows by lifting a loop through $[m]$:
its endpoint differs from its starting point by its
compact homology class divided by $m$ in the period
lattice.  The lift closes exactly when that homology
class is divisible by $m$.
On the compact primitive lattices the $[2]$ map is
$2I$, and the $[3]$ map is $3I$; these covering maps
are distinct from the ordinary maps which fill punctures.

Write $N=\ker(\Pi\to\mathcal Q)$ and regard it as a
subgroup of $K_2\leq\Sigma$ by (ASR.17).
The subgroup $K_2$ is characteristic in $\Sigma$, and
$N$ is characteristic in $K_2$; hence $N$ is
characteristic in $\Sigma$.
The quotient $\mathcal Q$ retains compact homology
modulo three.  Since the actual lattice map of $[2]$
is $2I$, which is invertible modulo three, this gives
$N\subseteq K_2\cap K_3=K_6$.
Define the ordinary characteristic quotient and its
three named subgroups by
\[
 \mathcal P=\Sigma/N,\qquad
 \mathcal P_2=K_2/N=\mathcal Q,\quad
 \mathcal P_3=K_3/N,\quad
 \mathcal P_6=K_6/N.                                   \tag{ASR.18}
\]
Their indices in $\mathcal P$ are respectively
$4,9,36$.  In particular $|\mathcal P|=4|\mathcal Q|$.
Every inclusion and quotient in (ASR.18) is induced
by the same subgroup $N$ of the same based free group.

For completeness the actual common geometric square is
\[
\begin{array}{ccc}
 E\setminus E[6]&\xrightarrow{[2]}&E\setminus E[3]\\
 {\scriptstyle[3]}\downarrow&&\downarrow{\scriptstyle[3]}\\
 E\setminus E[2]&\xrightarrow{[2]}&E\setminus\{0\}.
\end{array}                                             \tag{ASR.19}
\]
It is Cartesian.  A pair $(u,v)$ in the fibre product
satisfies $3u=2v\ne0$, and its inverse image is
$w=v-u$: then $2w=u$, $3w=v$, and $6w\ne0$.
Choose $w_0\notin E[6]$ and use the basepoints
$w_0,2w_0,3w_0,6w_0$ at the four respective corners.
All arrows in (ASR.19) are then literal based covering
maps.  On fundamental groups, it is precisely the
inclusion square $K_6\subset K_3,K_2\subset\Sigma$,
and its finite quotient is the square in (ASR.18).
The ordinary inclusions which fill the extra punctures,
such as $E\setminus E[6]\hookrightarrow E\setminus E[3]$,
are different named maps.  Remembering and then forgetting
these marked subsets gives the ordinary mapping-class
maps $\Gamma_{1,36}\to\Gamma_{1,9}$ and
$\Gamma_{1,36}\to\Gamma_{1,4}$.  They do not replace
the multiplication maps in (ASR.19).

It remains to check that passing to the ordinary outer
action on $\mathcal P$ has not introduced a deck kernel
in (ASR.2).  Use a representative of a modular mapping
class fixing the based peripheral word $c$ around the
origin in $\Sigma$.  Its chosen lift to $K_2$ fixes the
origin peripheral word in (ASR.8), which is the lift
of that same $c$.
An inner action on $\mathcal P$, say by $v$, therefore
satisfies $vcv^{-1}=c$ in $\mathcal P$.
Its image in $\mathcal P/\mathcal Q=C_2^2$ is the
deck translation acting on the four punctures of
$E\setminus E[2]$.
Their homology classes over $\mathbb F_3$ are distinct:
three form the usual peripheral basis and the fourth
is their negative sum.  These classes are all retained
in $\mathcal Q$ by (ASR.15).
The equality $vcv^{-1}=c$ consequently forces that
deck translation to fix the origin label, and hence
to be the identity.  Thus $v\in\mathcal Q$.
The restricted action is inner on $\mathcal Q$, so
(ASR.2) gives the identity in $D^+$.
We have proved the additional actual finite surjection
\[
 \operatorname{im}\bigl(
 \operatorname{SL}_2(\mathbb Z)\to
                  \operatorname{Out}(\mathcal P)\bigr)
       \twoheadrightarrow D^+,                        \tag{ASR.20}
\]
with the same generator images (ASR.16).
This completes the ordinary surface realization and
the exact two-/three-division marking comparison.
\end{proof}

\par\medskip\noindent\textbf{Exact centralizer consequence and remaining comparison.}
Let $c=ab$, so that its ordinary signed matrix is $UL$.
The image in $\operatorname{SL}_2(\mathbb Z/12)$ of
$C_J(\rho(c))$ is contained in the actual $C_6$
generated by $UL$, by (ASR.2) and the complete finite
centralizer calculation of A.4.6.
The reverse containment is supplied by the six powers of
$\rho(c)$ themselves.  This is therefore an actual
surface realization of that decisive finite centralizer
image, with all six signed residues retained.

The remaining finite recognition problem is now located
at an actual common quotient: the deck group
$\mathcal P/\mathcal P_3=C_3^2$ acts on the specified
index-nine subgroup $\mathcal P_3$ in (ASR.18).
The four-puncture and thirty-six-puncture comparison maps
are the other two subgroups in that same diagram.
A.4.9 proves geometric comparison at its specified
class-two quotient; it has not been substituted for
geometric straightening at $\mathcal P_3$.
One may form their actual common refinement by
intersecting $N$ with the characteristic kernel of
(TPL.1) in $K_3$, both regarded inside $\Sigma$.
This retains both established finite components, but
the additional geometric straightening and the
complete framed five-point census at that refinement
are still to be proved.
